\subsection{Symmetric formal power series}

Let I be a set, \(\mathbf{X} = (\mathbf{X}_i)_{i\in I}\) a family of indeterminates and \(\mathbf{A}[[\mathbf{X}]]\) the algebra of formal power series in the \(\mathbf{X}_i\) . For every permutation \(\sigma \in \mathfrak{S}_I\) there exists a unique continuous automorphism \(\varphi_{\sigma}\) of the algebra \(\mathbf{A}[[\mathbf{X}]]\) which maps \(\mathbf{X}_i\) to \(\mathbf{X}_{\sigma (i)}\) for each \(i\in I\) (IV, p.~28, Prop. 4); it is clear that \(\sigma \mapsto \varphi_{\sigma}\) is a homomorphism of \(\mathfrak{S}_I\) into the group of continuous automorphisms of the algebra \(\mathbf{A}[[\mathbf{X}]]\) . Let \(f\in \mathbf{A}[[\mathbf{X}]]\) be a formal power series; we put \(\sigma f = \varphi_{\sigma}(f)\) and we shall say that the formal power series \(f\) is symmetric if \(\sigma f = f\) for each \(\sigma \in \mathfrak{S}_I\) . The symmetric formal power series form a closed subalgebra of \(\mathbf{A}[[\mathbf{X}]]\) which is denoted by \(\mathbf{A}[[\mathbf{X}]]^{\mathrm{sym}}\) and is equipped with the topology induced by that of \(\mathbf{A}[[\mathbf{X}]]\) .

Let T be an indeterminate. In the ring of formal power series A{[}{[}X, T{]}{]} the family \((\mathbf{X}_i\mathbf{T})_{i\in I}\) is summable, hence the family \((1 + \mathbf{X}_i\mathbf{T})_{i\in I}\) is multipliable (IV, p.~27, Prop. 2); moreover we have

\[
\prod_ {i \in I} (1 + X _ {i} T) = 1 + \sum_ {k \geqslant 1} s _ {k} T ^ {k},\tag{14}
\]

where the formal power series \(s_k \in \mathbf{A}[[\mathbf{X}]]\) is defined by

\[
s _ {k} = \sum_ {\mathrm{H} \in \mathfrak {P} _ {k}} \left(\prod_ {i \in \mathrm{H}} \mathrm{X} _ {i}\right) \quad (k \geqslant 1)\tag{15}
\]

(we denote by \(\mathfrak{P}_k\) the set of all \(k\) -element subsets of I). In particular we have \(s_1 = \sum_{i \in I} X_i\) . When I is finite with \(n\) elements we have \(s_k = 0\) for \(k > n\) ; more precisely, when \(\mathbf{I} = \{1,\dots,n\}\) then the formal power series \(s_k\) is nothing other than the elementary symmetric polynomial of degree \(k\) in \(\mathbf{X}_1,\ldots,\mathbf{X}_n\) .

Let \(\mathbf{S} = (\mathbf{S}_k)_{k\geqslant 1}\) be a sequence of indeterminates. Since the formal power series \(s_k\) is of order \(\geqslant k\) , and belongs to \(\mathbf{A}[[\mathbf{X}]]^{\mathrm{sym}}\) , conditions a) and b) of Prop. 4 of IV, p.~28 are satisfied with \(\mathrm{E} = \mathbf{A}[[\mathbf{X}]]^{\mathrm{sym}}\) ; there exists thus a unique continuous A-algebra homomorphism

\[
\varphi_ {I}: \mathrm{A} [ [ \mathrm{S} ] ] \rightarrow \mathrm{A} [ [ \mathrm{X} ] ] ^ {\text { sym }}
\]

such that \(\varphi_{I}(S_{k}) = s_{k}\) for each integer \(k\geqslant 1\)

THEOREM 2. --- a) If I is a finite set of n elements, then \(\varphi_{I}\) induces a bicontinuous isomorphism of A{[}{[}S \(_{1}\) , \ldots, S \(_{n}\) {]}{]} onto A{[}{[}X{]}{]} \(^{\text{sym}}\) .

\begin{enumerate}
\def\labelenumi{\alph{enumi})}
\setcounter{enumi}{1}
\tightlist
\item
  If I is infinite, \(\varphi_{\mathrm{I}}\) is a bicontinuous isomorphism of A{[}{[}S{]}{]} onto A{[}{[}X{]}{]} \(^{\mathrm{sym}}\) .
\end{enumerate}

In case a) we put \(\mathbf{B} = \mathbf{A}[[\mathbf{S}_1, \ldots, \mathbf{S}_n]]\) and equip this algebra with the topology induced by that of \(\mathbf{A}[[\mathbf{S}]]\) ; we also denote by \(\psi_{\mathrm{I}}\) the restriction of \(\varphi_{\mathrm{I}}\) to \(\mathbf{B}\) . In case b) we put \(\mathbf{B} = \mathbf{A}[[\mathbf{S}]]\) and \(\psi_{\mathrm{I}} = \varphi_{\mathrm{I}}\) . We shall equip the polynomial algebra \(\mathbf{A}[\mathbf{S}]\) with the graduation of type \(\mathbf{N}\) for which \(\mathbf{S}_k\) is of weight \(k\) for every integer \(k \geqslant 1\) .

Lemma 1. --- Let J be a finite subset of I, r an integer such that Card \(J \geqslant r\) and f a symmetric formal power series which is homogeneous of degree r in the \(X_{i}\) ( \(i \in I\) ). Let \(\bar{f}\) be the formal power series obtained by substituting 0 for \(X_{i}\) for each i in I - J. If \(\bar{f} = 0\) , then we have f = 0.

Let us put \(f = \sum_{|\alpha| = r} a_{\alpha} X^{\alpha}\) (where \(|\alpha|\) is the length \(\sum_{i \in I} \alpha_i\) of the multiindex \(\alpha = (\alpha_i)_{i \in I}\) ). If \(\alpha\) is a multiindex of length \(r\) , and \(J'\) is the support of \(\alpha\) (the set of \(i \in I\) such that \(\alpha_i \neq 0\) ), then Card \(J' \leqslant r\) , hence there exists a permutation \(\sigma \in \mathfrak{S}_I\) such that \(\sigma(J') \subset J\) . Put \(\beta_i = \alpha_{\sigma^{-1}(i)}\) for \(i \in I\) , then the monomial \(X^\beta = \prod_{i \in I} X_{\sigma(i)}^{\alpha_i}\) depends only on the indeterminates \(X_j (j \in J)\) , whence \(a_\beta = 0\) by the hypothesis \(\bar{f}=0\) . Since f is symmetric, we have \(a_{\alpha}=a_{\beta}\) and since \(\alpha\) was arbitrary, we conclude that f=0.

Lemma 2. --- Let \(f\) be a symmetric formal power series of degree \(r\) in the \(X_i (i \in I)\) . There exists a unique polynomial \(P \in B \cap A[S]\) , isobaric of weight \(r\) , such that \(f = \psi_I(P)\) .

The case when I is finite follows from Th. 1 (IV, p.~62).

Assume that I is infinite and choose a finite subset J of I, containing r elements. We keep the notation of Lemma 1; we remark that every isobaric polynomial of weight r in the \(S_{n}\) ( \(n \geqslant 1\) ) depends only on \(S_{1}, \ldots, S_{r}\) and that \(\bar{s}_{1}, \ldots, \bar{s}_{r}\) are the elementary symmetric polynomials in the r indeterminates \(X_{j}\) ( \(j \in J\) ). If P is an isobaric polynomial of weight r in the \(S_{n}\) and \(h = f - \psi_{\mathrm{I}}(\mathrm{P})\) , then \(\bar{h} = \bar{f} - \mathrm{P}(\bar{s}_{1}, \ldots, \bar{s}_{r})\) and Lemma 1 shows the relation \(f = \psi_{\mathrm{I}}(\mathrm{P})\) to be equivalent to \(\bar{f} = \mathrm{P}(\bar{s}_{1}, \ldots, \bar{s}_{r})\) . By Th. 1 (IV, p.~62) there exists a unique polynomial \(P \in A[S]\) , isobaric of weight r and such that \(\bar{f} = \mathrm{P}(\bar{s}_{1}, \ldots, \bar{s}_{r})\) , whence the result.

Lemma 3. --- For every integer \(m \geqslant 0\) let \(\mathfrak{c}_m\) be the ideal of the algebra \(\mathrm{A}[[\mathrm{X}]]^{\mathrm{sym}}\) consisting of all symmetric formal power series of order \(\geqslant m\) . The sequence \((\mathfrak{c}_m)_{m \geqslant 0}\) is a fundamental system of neighbourhoods of 0 in \(\mathrm{A}[[\mathrm{X}]]^{\mathrm{sym}}\) .

The lemma clearly holds when I is finite, so let us assume that I is infinite. For every finite subset J of I, consisting of m elements, denote by \(\tilde{J}\) the set of elements of \(\mathbf{N}^{(1)}\) of length \textless{} m and with support in J. Further denote by \(a_{j}^{\prime}\) the set of formal power series containing no term \(aX^{\alpha}\) with \(\alpha \in \tilde{J}\) . Since \(\tilde{J}\) is a finite subset of \(\mathbf{N}^{(1)}\) and every finite subset of \(\mathbf{N}^{(1)}\) is contained in a set of the form \(\tilde{J}\) , the family \((\mathfrak{a}_{\mathrm{j}}^{\prime})\) is a base of neighbourhoods of 0 in A{[}{[}X{]}{]} (IV, p.~26). Now Lemma 1 implies the relation \(a_{j}^{\prime} \cap A[[X]]^{\text{sym}} = c_{m}\) for every subset J with m elements, and this proves Lemma 3.

Since there are only a finite number of monomials of a given weight in the \(S_{k}\) , every formal power series \(f \in B\) may be written in a unique way as \(f = \sum_{r \geqslant 0} P_{r}\) , where \(P_{r}\) is an isobaric polynomial of weight r in \(B \cap A[S]\) . For every integer \(m \geqslant 0\) let \(\mathfrak{b}_m\) be the ideal of B consisting of all formal power series of the above type such that \(P_r = 0\) for \(0 \leqslant r < m\) . The sequence \((\mathfrak{b}_m)_{m \geqslant 0}\) is a base of neighbourhoods of 0 in B.

With the above notation \(\psi_{\mathrm{I}}(\mathbf{P}_{r})\) is a symmetric formal power series in the \(\mathbf{X}_i\) , homogeneous of degree \(r\) , and so this is the homogeneous component of degree \(r\) of \(\psi_{\mathrm{I}}(f)\) . Lemma 2 shows that \(\psi_{\mathrm{I}}\) is an algebra homomorphism of B onto \(\mathbf{A}[[\mathbf{X}]]^{\mathrm{sym}}\) , mapping \(\mathfrak{b}_m\) to \(\mathfrak{c}_m\) for every integer \(m \geqslant 0\) ; now Lemma 3 shows \(\psi_{\mathrm{I}}\) to be bicontinuous, and this completes the proof of Theorem 2.

